\section{Introduction}

Digital libraries increasingly steward research objects that live beyond the traditional narrative article \cite{wilkinson2016fair,smith2016softwarecitation,barker2022fair4rs,soilandreyes2022rocrate}. Datasets, software, scripts, model weights, evaluation assets, and external repositories have become fundamental components of the scholarly record, ensuring that computational findings remain inspectable, reusable, and preservable. Open \llm research intensifies this paradigm shift, as a single paper may release a small model, a LoRA adapter, a quantized checkpoint, a merged model, a distilled student, or a deployment package whose operational meaning depends heavily on upstream dependencies and visible release assets. These compact and derived artifacts are rarely contained within a single monolithic record, existing instead as scattered traces across article text, \hf repositories, model cards, source code links, platform metadata fields, family names, and release statements. Institutional curation must therefore address a distributed scholarly record rather than an isolated paper or repository.

This distribution creates a practical curation problem because public visibility does not automatically guarantee curation readiness. The artifact is not necessarily hidden; it may be visible in several places at once. Yet a future curator may still be unable to decide what object should be cited, which repository is the release rather than a dependency, which upstream relation should be recorded, or which assets should be selected for preservation triage. A visible model-hub link may help a researcher locate an asset, but it rarely clarifies whether the target repository represents an official release, a baseline dependency, a temporary checkpoint package, or a third-party derivative. Similarly, a referenced model family name can identify the target artifact, an upstream base checkpoint, a teacher model, an experimental baseline, or general background context. In this sense, the risk is not only artifact loss but evidence loss: public traces exist, while the relations that make them curatable remain underspecified. The artifact may remain online, but the evidence needed to cite, attribute, trace, reuse, and preserve it can still disappear into untyped links and ambiguous release statements.

Existing literature provides critical foundational insights without offering a collection-scale assessment of this systemic record problem. Prior work on model cards, datasheets, model openness, reproducibility, software citation, data citation, research-object packaging, and context-based URL classification demonstrates that computational artifacts require structured documentation rather than binary availability metrics \cite{gebru2021datasheets,mitchell2019modelcards,pineau2021reproducibility,bommasani2023fmi,white2024mof,salsabil2025context}. Empirical studies of model hubs and asset reuse further clarify how these models circulate through open-source artificial intelligence communities \cite{jiang2023ptmreuse,taraghi2024modelreuse,osborne2024aicommunity}. Most of these investigations, however, evaluate only one record type at a time, focusing exclusively on isolated papers, individual URLs, code repositories, distinct model cards, datasets, or platform-specific model-hub behavior. Digital preservation systems still lack empirical evidence demonstrating how papers, model hubs, model cards, and release links interact to support or hinder the systematic curation of compact and derived open \llm artifacts.

We address this empirical gap by operationalizing \curatability as a record-level measure of curation readiness. In this study, a curatable public record provides sufficient evidence to determine the artifact's operational role, its record-supported scholarly linkage, its upstream model evidence, and its visible release assets available for inspection and preservation triage. Utilizing a fixed May 2026 snapshot, we evaluate 191,375 public \hf model repositories and 2,729 scholarly paper records, filtering the literature into a core analysis set of 2,214 papers. Our measurement pipeline constructs linked repository-paper evidence by normalizing repository metadata, paper-side links, alignment signals, provenance evidence, and release-package fields, incorporating expert audits and sensitivity checks to calibrate boundary conditions.

Cross-system measurements reveal a severe drop-off from initial asset visibility to ultimate curation readiness. Public records successfully establish a broad discovery layer, with 90.7\% of analysis-set papers containing at least one useful curation signal. The proportion of viable records contracts dramatically when curation requires multiple evidence fields to co-occur. Explicitly matched paper-to-\hf repository evidence covers only 5.6\% of the analysis-set papers. An operational threshold requiring usable provenance alongside concrete release evidence identifies 18.1\% of records, whereas a stricter threshold requiring usable provenance, asset-level release evidence, and a direct linkage signal identifies only 6.1\% of the corpus. Systemic fragmentation across metadata platforms, rather than an absolute absence of public artifacts, forms the primary bottleneck to preservation-aware digital-library curation.

Primary contributions of this work encompass three distinct dimensions:

\begin{itemize}[leftmargin=*]
    \item \textbf{Digital-library framing.} We establish a curation framing that repositions compact and derived open \llm artifacts as distributed scholarly objects.
    \item \textbf{Audited empirical record construction.} We present a linked record construction for evaluating artifact identity, scholarly linkage, upstream evidence, release-package evidence, and integrated curation readiness.
    \item \textbf{Metadata and interface recommendations.} We translate these empirical insights into concrete interventions for model hubs, scholarly indexes, and digital libraries, detailing fields for typed artifact roles, bidirectional paper--repository relations, upstream relation types, provenance evidence tiers, and structured release-package fields.
\end{itemize}

These findings offer a constructive implication for open science, demonstrating that current infrastructure supports broad public visibility, while the next milestone requires rendering public evidence sufficiently structured to support systematic citation, reuse, provenance tracking, and preservation triage.
